\begin{reviewercomment}
\textbf{5. Evaluation via Proper Scoring Rules:} The predictive evaluation relies mainly on MAE and RMSE (point prediction metrics). Given the emphasis on reliable uncertainty quantification, report proper scoring rules, such as Continuous Ranked Probability Score (CRPS) and interval score, for models providing predictive distributions.
\end{reviewercomment}
\begin{response}
We thank the reviewer for this comment. We added two proper scoring rules, the Continuous Ranked Probability Score (CRPS) and the interval score, alongside RMSE and MAE for all five models (Gaussian process with and without features, Bayesian ridge, MLP, XGBoost), for both Nigeria and Chad.

For an observation $y$ and a predictive cumulative distribution function $F$, the CRPS is
\begin{equation}
  \mathrm{CRPS}(F, y) = \int_{-\infty}^{\infty} \big(F(x) - \mathbb{1}\{x \ge y\}\big)^2 \, dx,
\end{equation}
a strictly proper scoring rule that reduces to the absolute error $|y - x|$ when $F$ is a point mass at $x$, so it generalises MAE to full predictive distributions. Lower values indicate better calibrated and sharper forecasts. For each model we obtain $M = 100$ draws $x_1, \dots, x_M$ from its predictive distribution (posterior draws for the Gaussian process, MC dropout draws for the MLP, bootstrap predictions for XGBoost, and Monte Carlo draws from the fitted predictive Normal, on the logit scale, back-transformed to the FCS scale, for Bayesian ridge) and compute the empirical estimator
\begin{equation}
  \widehat{\mathrm{CRPS}}(x_{1:M}, y) = \frac{1}{M} \sum_{m=1}^{M} |x_m - y| - \frac{1}{2M^2} \sum_{m=1}^{M} \sum_{m'=1}^{M} |x_m - x_{m'}|.
\end{equation}

For a $(1-\alpha)$ central prediction interval $[\ell, u]$, the interval score is
\begin{equation}
  \mathrm{IS}_\alpha(\ell, u, y) = (u - \ell) + \frac{2}{\alpha}(\ell - y)\,\mathbb{1}\{y < \ell\} + \frac{2}{\alpha}(y - u)\,\mathbb{1}\{y > u\},
\end{equation}
which rewards narrow intervals and penalises observations falling outside them, with the penalty inversely proportional to $\alpha$ \cite{gneiting2007strictly}. We report $\mathrm{IS}_{0.05}$, matching the 95\% intervals already used for the coverage and interval-length diagnostics. Both scores are computed per observation and averaged across test folds and regions.

Tables~\ref{tcd-tab-models} and \ref{nga-tab-models} report the updated results. The Gaussian process with features achieves the lowest CRPS for both countries (0.029 for Chad, 0.032 for Nigeria), consistent with its already lowest RMSE and MAE, confirming that it produces the best calibrated and sharpest predictive distributions overall. The interval score reveals a limitation not visible from RMSE or MAE alone: XGBoost obtains a competitive RMSE (0.054 for Chad, 0.064 for Nigeria) but by far the worst interval score (1.25 and 1.34), because its 95\% intervals are severely overconfident (coverage of 23\% and 28\%, against a nominal 95\%). The interval score therefore penalises XGBoost's narrow but unreliable intervals, which point metrics alone would not reveal.
\end{response}
\begin{change}
Section X, page Y, lines Z--W.
\end{change}
